\section{Exact ordinary density for imaginary quadratic fourfolds}
\label{sec:fourfold}

The preceding bounds allow a base extension. For a fourfold with full
geometric endomorphism algebra imaginary quadratic, the Hodge signature
determines the ordinary density over the original field together with the
field of definition of the endomorphisms.

\begin{theorem}\label{thm:fourfold-density}
Let \(A/K\) be an abelian fourfold over a number field with
\(E=\operatorname{End}^0(A_{\overline K})\) imaginary quadratic. Let
\(K_E/K\) be the field defining the geometric endomorphisms. The natural
density of ordinary primes of \(A/K\) is
\[
 \begin{cases}
  1 & \text{in signature }(2,2),\\
  [K_E:K]^{-1} & \text{in unordered signature }(1,3).
 \end{cases}
\]
These are the only signatures under the stated hypothesis.
\end{theorem}

If the signature were \((0,4)\), every \(E\)-linear endomorphism of
rational Betti cohomology would preserve its Hodge decomposition,
because each coefficient eigenspace would have a single Hodge type.
The geometric endomorphism algebra would then contain \(M_4(E)\).
In signature \((1,3)\), the hypothesis of simple signature, relative
rank four and the Galois property of \(E/\Q\) give the density by
\cite[Theorem~4.18]{cantoral}. To identify \(K_E\), write \(\tau\)
for the unique embedding of Hodge multiplicity one. If a Galois
element \(g\) acts on \(E\) by \(\gamma_g\), it sends that Hodge
eigenspace to the one labelled \(g\circ\tau\circ\gamma_g^{-1}\).
Uniqueness gives \(g\circ\tau=\tau\circ\gamma_g\), hence \(K_E=K\,\tau(E)\).

In the balanced case, Fit\'e proves positive ordinary density over
\(K\) for the wider class allowing an imaginary quadratic subfield of
the geometric endomorphism algebra. His proof gives density one after
finite extension \cite[Theorem~4 and \S8]{fite}. To obtain density one
over \(K\) with full endomorphism algebra \(E\), we treat the components
where Frobenius conjugates \(E\) directly. On these components Frobenius
interchanges two symplectically dual spaces. A normalized determinant
of its square then has even valuation, which turns a nonordinary
reduction into a bounded integer condition on Frobenius.

\subsection{Monodromy and the middle coefficient}

Assume signature \((2,2)\). Fix a polarization over \(K\). Its Rosati
involution acts on \(E\) by complex conjugation, the positive
involution of an imaginary quadratic field. Fix an auxiliary prime
\(\ell\) split in \(E\). Let \(G\) be the algebraic monodromy group
on first cohomology and let \(D=(G^\circ)^{\mathrm{der}}\). Over
\(\overline{\Q}_\ell\) write the two \(E\)-eigenspaces as \(V_+,V_-\).
The polarization pairs them perfectly. After a finite extension
defining the endomorphisms and making monodromy connected, Faltings's
theorems \cite[Theorems~3--4]{faltings} identify the commutant of
\(G^\circ\) with \(E\otimes\overline{\Q}_\ell\) and give semisimplicity.
Thus \(V_+,V_-\) are irreducible and inequivalent under \(G^\circ\).

The two relative determinant characters equal \(\chi^2\) on
\(G^\circ\), where \(\chi\) denotes the similitude. Indeed, after
extending the base to contain \(E\) and define its action, the
determinants have infinity type \((2,2)\) by
\cite[Proposition~14]{fite}. Division by \(\chi^2\) gives an algebraic
Hecke character of trivial infinity type and hence finite image. For
the scalar characters \(z_+,z_-\) of the connected center, the
determinant and polarization give \(z_+^4=\chi^2\) and
\(z_+z_-=\chi\). Thus \(z_+^2=\chi\) by connectedness, and
\(z_+=z_-\). The representation of \(D\) on \(V_+\) is irreducible
because the connected center acts by scalars. Self-duality would make
the two \(G^\circ\)-modules isomorphic, since the central characters
agree. The classification of four-dimensional irreducible semisimple
modules gives \(\mathrm{SL}_4\) on \(V_+\) \cite[\S3.1]{hui4}. The
three other representations in that list are self-dual. In dual bases
the action of \(D\) on \(V_+\oplus V_-\) is
therefore exactly \(\operatorname{diag}(S,S^{-\mathsf T})\) for
\(S\in\mathrm{SL}_4\). The finite extensions in this identification
leave \(G^\circ\) and its representations unchanged.

Now work over the original \(K\). At a good absolute degree-one prime
of residue characteristic \(p\), omit \(p=\ell\), primes ramified in
\(E\) or \(K_E\), and primes dividing the endomorphism order index or
polarization degree. Write \(\pi\) for the \(p\)-power Frobenius of the
reduction. Its similitude is \(p\). On each monodromy component the
action on \(E\) is constant, and a representative either preserves the
two blocks or interchanges them.

Suppose first that \(\pi\) preserves the blocks. Its relative second
coefficient \(a_2\) lies in \(\mathcal O_E\) by the trace comparison in
Section~\ref{sec:coefficients}. The geometric coefficient blocks have
slope sum two, because both Hodge multiplicities are two. At an inert
prime this follows by taking the determinant around the two-step
Dieudonn\'e cycle and dividing its valuation by two. If the reduction
is nonordinary, both slope lists differ from \((0,0,1,1)\), by
polarization. A nonordinary block has a positive sum of its two
smallest slopes, since all slopes lie in \([0,1]\) and sum to two.
Its second coefficient therefore has positive integral valuation over
the unramified coefficient field. The residue degrees sum to two,
so \(p^2\) divides \(\Nm_{E/\Q}(a_2)\). The Weil bound then makes
\(\Nm_{E/\Q}(a_2)/p^2\) an integer of absolute value at most \(36\).

This quotient is the Frobenius value of the regular class function
\(a_{2,+}(g)a_{2,-}(g)/\chi(g)^2\) on the components preserving the
blocks. For a matrix \(A_+\) of \(g\) on \(V_+\), its matrix on
\(V_-\) is \(\chi(g)A_+^{-\mathsf T}\), so the function is
\(a_2(A_+)^2/\det A_+\). Multiplication by \(D\) varies \(A_+\)
over all matrices of its determinant. In particular, at matrices
\(a\operatorname{diag}(u,u^{-1},1,1)\) with \(a^4=\det A_+\), it
takes the values \((2+2u+2u^{-1})^2\). It is nonconstant on every
such component.

\subsection{Pfaffian parity on the interchanging components}

\begin{proposition}\label{prop:fourfold-pfaffian}
At a prime as above where \(\pi\) interchanges the blocks, a
nonordinary reduction makes
\[
  \Nm_{E/\Q}\!\left(\det_E(1+\pi^2/p)\right)
\]
an integer of absolute value at most \(256\).
\end{proposition}

\begin{proof}
Extend the crystalline lattice to a finite unramified Witt ring
\(W\) that splits the coefficient field and defines the specialized
action. The lattice has two rank-four summands \(M_+,M_-\) in perfect
duality. The Frobenius cokernel is killed by \(p\) and the corresponding
Hodge quotient has dimension two at each label
\cite[\S\S1.1.3 and 1.3.1]{moonencanonical}. Thus the semilinear
Frobenius on the Dieudonn\'e lattice has elementary divisors
\(1,1,p,p\) on each map between coefficient summands. Over \(W\) it
is the linear pullback \(\pi\) together with Witt Frobenius on scalars.
The scalar map permutes the coefficient summands by lattice
isomorphisms. Thus the block maps of linear \(\pi\) also have
determinant valuation two.

In dual integral bases write
\[
 \pi=\begin{pmatrix}0&C\\ B&0\end{pmatrix}.
\]
The similitude condition gives \(C=-p B^{-\mathsf T}\). On \(M_+\)
we obtain
\[
 \det(1+\pi^2/p)=\frac{\det(B^{\mathsf T}-B)}{\det B}
               =\frac{\operatorname{Pf}(B^{\mathsf T}-B)^2}{\det B}.
\]
Its valuation is even or infinite. The two blocks of \(\pi^2\) are
conjugate by \(\pi\), and their slopes are complementary under the
polarization. Thus the valuations of the eigenvalues of \(\pi^2/p\)
in each block have the form \((-a,-b,b,a)\) with
\(0\le b\le a\le1\). If the reduction is nonordinary then \(a+b<2\).
The determinant on each block has valuation at least \(-a-b>-2\),
by applying \(v_p(1+x)\ge\min(0,v_p(x))\) to each eigenvalue. By
parity it is integral at \(p\).

The same trace comparison as for \(\pi\) applies to the endomorphism
\(\pi^2\), which commutes with \(E\). The relative characteristic
coefficients are integral and agree in the Tate and crystalline
realizations. The displayed norm in the proposition is rational
with denominator a power of \(p\). It has just been shown integral
at \(p\). Its complex bound is \(2^8\), since all eigenvalues of
\(\pi^2/p\) have complex modulus one.
\end{proof}

The integer in Proposition~\ref{prop:fourfold-pfaffian} is the
Frobenius value of the regular class function
\[
 g\longmapsto
 \prod_{\sigma\in\{+,-\}}
   \det\!\left(1+g^2|_{V_\sigma}/\chi(g)\right).
\]
On any interchanging component, multiplication by \(D\) varies its
off-diagonal block \(B\) over all matrices of its determinant. Choose
\(b\) with \(b^4=\det B\). At \(B=b I_4\) the determinant ratio in
the proof is zero. At \(B=b J_4\), with \(J_4\) the standard
alternating matrix of determinant one, it is \(2^4\). The two
block determinants agree because \(g\) interchanges the blocks and
commutes with \(g^2\). The class function takes the values zero and
\(2^8\) on each such component, and is nonconstant.

The two functions used above descend to \(\Q_\ell\), since their
products are invariant under permuting the two embeddings of \(E\).
Each of their finitely many integer levels is a proper closed subset
on every component on which the function was used. Algebraic
Chebotarev, in the form used in Section~\ref{sec:adjoint}, makes their
union a density-zero condition on primes. Absolute degree-one primes
have density one, so this completes the balanced case of
Theorem~\ref{thm:fourfold-density} over \(K\).

\subsection{An explicit genus-four Jacobian}

\begin{example}\label{ex:fourfold-curve}
Let \(C/\Q\) be the smooth projective curve with affine equation
\[
 y^3=A(x)B(x)^2,
 \qquad A(x)=x^3+x+1,\qquad B(x)=x^3-2x^2-2x+1.
\]
Its Jacobian has geometric endomorphism ring \(\Z[\zeta_3]\), with
endomorphism field \(\Q(\zeta_3)\). Its ordinary primes over \(\Q\)
have natural density one.
\end{example}

The discriminants of \(A,B\) are \(-31,125\), and their resultant is
\(31\). The six finite branch points are distinct and infinity is
unramified. Riemann--Hurwitz gives genus four, and the same tame
local models give good reduction away from \(3,5,31\). The deck
action by \(y\mapsto\zeta_3y\) supplies \(E=\Q(\zeta_3)\), as its
invariant cohomology is that of \(\mathbb P^1\). The forms
\(x^j\,dx/y\) and \(x^j B(x)\,dx/y^2\), for \(j=0,1\), are
holomorphic and give two independent vectors in each deck
eigenspace. At the branch points this follows from the orders of
\(dx\) and \(y\), which are two and either one or two, respectively.
At infinity \(x\) and \(y\) have pole orders one and three. The
signature is therefore \((2,2)\).

Here is the finite calculation certifying the endomorphisms. Put
\(\zeta=\zeta_3\). At \(p=7,13\) use the primitive cube roots
\(2,3\) in \(\mathbb F_p\), respectively. For \(q=p^r\) with
\(1\le r\le4\), count \(n_i\) values of \(x\) at which
\(A(x)B(x)^2\ne0\) and its \((q-1)/3\)-th power is the \(i\)-th
power of the chosen cube root. The \(j\)-th deck power composed with
\(q\)-Frobenius fixes all three points above each such \(x\) with
index \(-j\) modulo three. Its fixed branch points contribute the
same number for each \(j\), and the three points at infinity are
fixed for \(j=0\) alone. Thus Fourier projection of the Lefschetz
trace gives \(-1-\sum_{i=0}^2n_i\zeta^i\) on one coefficient block,
with the conjugate trace on the other. Newton identities applied
to the four power traces give the relative polynomials
\[
\begin{aligned}
 P_7(X)&=X^4+(2-\zeta)X^3+(3-3\zeta)X^2
                    +(7-14\zeta)X-49\zeta,\\
 P_{13}(X)&=X^4+(5+\zeta)X^3+16X^2+(52-13\zeta)X+169.
\end{aligned}
\]
The middle coefficients of \(P_p\overline{P_p}\) are \(167\) and
\(1101\), units at the respective primes. Both reductions are
ordinary and both relative determinants have valuation two. In
particular the four slopes in each relative block are \(0,0,1,1\).
The squarefree relative factorizations have the following degrees.
\[
\begin{array}{c|c|cc}
\ell&\text{image of }\zeta&P_7&P_{13}\\ \hline
73&8&1,3&4\\
79&23&4&1,3
\end{array}
\]
These give irreducibility over \(E\) and Galois group \(S_4\) for
each quartic. The three-cycle and four-cycle force its order to be
divisible by twelve, and the odd four-cycle rules out \(A_4\).
The two conjugate quartics are distinct at each \(p\). Thus the
absolute polynomial is irreducible and each root field \(L_p\)
contains \(E\), since \(L_p\) and \(L_pE\) both have degree eight
over \(\Q\). There is no proper intermediate field between \(E\)
and \(L_p\), by maximality of the point stabilizer in \(S_4\).

The eight roots at each \(p\) remain distinct after any positive
power. A ratio of distinct roots in one relative quartic has a
conjugate of nonzero valuation, since an \(S_4\) permutation sends
the two roots to different valuations. By polarization the conjugate
roots are \(p/\alpha_j\) for roots \(\alpha_j\) in the first quartic.
A ratio across the two quartics is then \(\alpha_i\alpha_j/p\), whose
valuation is nonzero after sending the selected roots to valuation
zero. Each power of Frobenius consequently generates the same degree
eight field \(L_p\). Tate's theorem gives endomorphism algebra \(L_p\)
over every finite extension of the residue field, and hence over its
algebraic closure \cite[Theorem~2(c)]{tatefinite}.

The two absolute polynomials are squarefree modulo two, with factor
degrees \((4,4)\) and \((2,6)\), respectively, so the two root fields
are nonisomorphic. Specialization injects the geometric endomorphism
algebra in characteristic zero into both fields, compatibly with
the \(E\)-actions. Since neither root field has a proper intermediate
field over \(E\), this proves that the algebra is \(E\). The deck
action supplies its maximal order and is defined over
\(\Q(\zeta_3)\), where the nontrivial Galois automorphism conjugates
the deck generator to its inverse. Theorem~\ref{thm:fourfold-density}
now applies. The exact calculation is reproduced in the supplement,
which also checks the ordinary reduction at the inert prime \(11\).
